% !TeX root = ../../Book3.tex
% 纲要标题：## 第16章　深度与 Serre 条件
% 纲要位置：工作记录/计划纲要.md，第 2485 行。
% 纲要细目（写作范围）：

\chapter{深度与 Serre 条件}
\label{b3:ch:16}
\BookChapterNumber{b3:ch:16}

\begin{chapterabstract}
深度把局部环上的有限生成模中能够依次选取的正则元素个数变成一个不依赖选择的整数。本章先用 Koszul 同调比较所有极大正则序列，再证明深度等于剩余域的 Ext 首次非零次数。借助这一刻画，建立短正合列中的深度引理、关联素理想的维数估计和局部化不等式，并把各点的深度要求整理为 Serre 条件。这些结果为投射维数、Cohen--Macaulay 模与局部上同调提供共同的计算依据。
\end{chapterabstract}

\begin{learninggoals}
\begin{itemize}
\item 证明极大正则序列同长度，区分“不能延长”与预先假定“长度最大”。
\item 说明有限性在深度零判据、Ext 模及局部化估计中的作用。
\item 用短正合列比较深度，并辨别何时只有下界、何时可以取得等式。
\item 在全部相关素理想处检查 \((S_k)\)，用关联素理想解释 \((S_1)\) 的含义。
\end{itemize}
\end{learninggoals}

设 \(k\) 为域。在局部环 \(R=k[x,y]_{(x,y)}/(xy)\) 中，以下用 \(x,y\) 表示它们的剩余类。\(x\) 和 \(y\) 都是零因子，但 \(x+y\) 不是：它在两个整环商 \(R/(x)\)、\(R/(y)\) 中的像都非零，而 \(R\) 嵌入这两个商的直和。因此仍可从极大理想中选出一个非零因子。若改为 \(k[x,y]_{(x,y)}/(x^2,xy)\)，非零的 \(\bar x\) 被极大理想的每个元素杀死，连第一步都无从开始。两个环的 Krull 维数都为一，正则元素可选取的长度却不同。维数测量素理想链的长度，深度则测量还能连续选取多少个非零因子；后一环的 \(y\) 方向仍然存在，但闭点处由 \(\bar x\) 生成的部分已阻止任何正则切割。

\ProseParagraphBreak

若只按某次选择数出正则元素，就要证明换一种选择不会改变所得最大长度。Koszul 同调与剩余域的 Ext 给出不依赖选择的检验方式；短正合列与局部化又让这种数值在模之间、在不同点之间比较。由此得到的深度不等式，将把可见的维数与隐藏的嵌入结构分开。

\ProseParagraphBreak

本章使用\chapref{b3:ch:15}的正则序列与 Koszul 同调，以及\chapref{b3:ch:04}的关联素理想、\chapref{b3:ch:09}的维数和参数系。\secref{b3:sec:1505}已证明自由消解的存在性，以及比较映射在同伦意义下的唯一性；\secref{b3:sec:1602}据此定义 Ext，建立双变量长正合列，并证明所需的有限性和局部化性质。第四卷第6—9章系统讨论这些同调函子的理论，\chapref{b3:ch:22}则用局部上同调给出深度的另一种刻画。

\ProseParagraphBreak
正则序列长度与 Ext、Koszul 同调之间的关系，可对照松村英之的《Commutative Ring Theory》\cite[\S16, Theorems 16.6--16.8, pp. 129--131]{Matsumura1989CommutativeRingTheory}。本章用于局部化估计的消失引理和关联素理想维数界，对应该书的 Theorems 17.1--17.2\cite[pp. 133--134]{Matsumura1989CommutativeRingTheory}；环的 Serre 条件在该书 \S23, p. 183 中集中说明\cite[\S23, p. 183]{Matsumura1989CommutativeRingTheory}。本章同时明确有限生成模版本所使用的支集维数。

% !TeX root = ../../Book3.tex
% 纲要标题：### 16.1 深度
% 纲要位置：工作记录/计划纲要.md，第 2487 行。
% 纲要细目（写作范围）：
% * 极大正则序列长度
% * 深度的有限性与存在性、基座及深度零与有限长度子模的判据
% * 深度不超过维数；非约化环的正深度

\section{深度}
\label{b3:sec:1601}

维数数的是素理想链，正则序列则问有多少个元素能够依次以单射作用。后一种长度似乎依赖选择：一个不能再延长的序列，是否可能比另一个短？局部有限性使这些长度相同，Koszul 同调给出一个可以同时比较所有选择的量，其共同长度就是深度。


% 纲要内容（仅作写作计划，尚非教材正文）：
%
% * 极大正则序列长度
% * 深度的有限性与存在性、基座及深度零判据
% * 深度不超过维数
%

\subsection{极大正则序列的长度}
\label{b3:subsec:1601:01}

\begin{definition}\label{b3:def:maximal-regular-sequence}
设 \((R,\mathfrak m)\) 为交换 Noether 局部环，\(M\ne0\) 为有限生成的 \(R\) 模。若 \(x_1,\ldots,x_t\in\mathfrak m\) 为 \(M\) 正则序列，而且不能再添上任何 \(y\in\mathfrak m\) 得到更长的正则序列，则称它为 \(\mathfrak m\) 中的一个\term[jida-zhengze-xulie]{极大正则序列}[maximal regular sequence]。这里“极大”指不能延长，尚未预先断言其长度最大。
\end{definition}

不能延长只说明当前选择已到尽头，一般并不能保证不同选择具有相同长度。下面用同一组极大理想生成元的 Koszul 同调比较这些选择，证明正则序列在这里确实有一个共同的最大长度。

\begin{theorem}\label{b3:thm:maximal-regular-length}
设 \((R,\mathfrak m)\) 为交换 Noether 局部环，\(M\ne0\) 为有限生成的 \(R\) 模。每个由 \(\mathfrak m\) 中元素构成的 \(M\) 正则序列都能延长为极大正则序列，且所有极大正则序列长度相同。若取 \(\mathfrak m=(y_1,\ldots,y_e)\) 的任意有限生成组，不必为极小生成组，其共同长度为
\[
e-\max\{\,i\mid H_i(y_1,\ldots,y_e;M)\ne0\,\}.
\]
\end{theorem}
\begin{proof}
每添加一个正则元素，已选元素生成的理想严格增大：若新元素在原理想中，它在非零的逐次商上作用为零，不可能单射。Noether 环的理想升链条件保证这个过程在有限步停止。

先看非零有限生成模 \(N\) 上的空序列何时已经极大。对 \(x\in\mathfrak m\)，Nakayama 引理保证 \(N/xN\ne0\)，所以能否添加 \(x\) 只取决于乘 \(x\) 是否单射。空序列极大便意味着 \(\mathfrak m\) 中每个元素都是 \(N\) 的零因子。由\cref{b3:thm:associated-zero-divisors}，
\[
\mathfrak m\subseteq\bigcup_{\mathfrak p\in\operatorname{Ass}_RN}\mathfrak p.
\]
因 \(R\) Noether 且 \(N\) 有限生成，\(\operatorname{Ass}_RN\) 是有限集，故\cref{b3:lem:finite-prime-avoidance}适用，给出 \(\mathfrak m\) 包含于某个 \(\mathfrak p\in\operatorname{Ass}_RN\)。局部环中每个素理想都包含于 \(\mathfrak m\)，于是 \(\mathfrak p=\mathfrak m\)。存在 \(n\ne0\) 满足 \(\mathfrak mn=0\)。因此 Koszul 最高同调
\[
H_e(y;N)=\{\,n\in N\mid\mathfrak mn=0\,\}
\]
非零，且没有次数超过 \(e\) 的项。反过来，这样的 \(n\) 也使 \(\mathfrak m\) 中任何元素都不可能正则。因此找不到第一个正则元素时，已有同一个非零元素被整个极大理想零化，而不是只有不同元素各自被某些零因子零化。

对任意非零有限生成模 \(N\)，令 \(q(N)=\max\{i\mid H_i(y;N)\ne0\}\)。这个最大值存在：由\cref{b3:thm:nakayama}，\(H_0=N/\mathfrak mN\ne0\)，而复形只有次数零到 \(e\) 的项。要比较取正则商前后的最高非零次数，关键是 \(x\in\mathfrak m=(y_1,\ldots,y_e)\)，所以\cref{b3:prop:koszul-signs-homotopy}给出 \(xH_i(y;N)=0\)。

若 \(x\in\mathfrak m\) 在 \(N\) 上正则，将短正合列 \(0\rightarrow N\xrightarrow{x}N\rightarrow N/xN\rightarrow0\) 张量 \(K_\bullet(y;R)\)。该复形的每一项都是有限秩自由模，故得到逐项短正合列。按\cref{b3:lem:tor-basic-properties}证明中的圈提升构造，它诱导同调长正合列。在 \(H_i(y;N/xN)\) 的两侧，长列的乘法 \(x\) 都为零，因此长列分成短正合列
\[
0\longrightarrow H_i(y;N)\longrightarrow H_i(y;N/xN)
\longrightarrow H_{i-1}(y;N)\longrightarrow0
\]
正合。令 \(q=q(N)\)，则 \(H_{q+1}(y;N)=0\)、\(H_q(y;N)\ne0\)。取 \(i=q+1\)，上列给出
\[
H_{q+1}(y;N/xN)\xrightarrow{\sim}H_q(y;N)\ne0.
\]
当 \(i>q+1\) 时，短正合列的两端均为零，中间项也为零。因此每除去一个正则元素，最高非零 Koszul 同调次数恰好上移一位，即 \(q(N/xN)=q(N)+1\)。若极大正则序列长为 \(t\)，它的非零末端商已找不到正则元素，前面的空序列判据说明该商的最高非零同调次数为 \(e\)。连续应用这个等式便得到 \(e=q(M)+t\)。右端除 \(t\) 外不依赖所选序列，故所有长度相同。
\end{proof}

这里的 \(e\) 是所选生成组的长度，不一定等于嵌入维数。例如，对域 \(k\)，在 \(R=k[X,Y]_{(X,Y)}\) 中，\(X,Y\) 是正则序列，所以 \(H_0(X,Y;R)=k\)，正次数同调为零，最高非零次数为 \(0\)。也可以用 \(X,Y,X\) 生成极大理想。把第三个元素减去第一个元素，由\cref{b3:prop:koszul-generator-change}将其复形换成 \(K_\bullet(X,Y,0;R)\)。张量分解中的 \(K_\bullet(0;R)\) 在次数零、一各有一个 \(R\)，微分为零，故
\[
H_i(X,Y,X;R)
\cong H_i(X,Y;R)\oplus H_{i-1}(X,Y;R).
\]
于是这一冗余生成组的零次和一次同调都是 \(k\)，更高同调为零；公式中的差由 \(2-0\) 变成 \(3-1\)，共同长度仍为 \(2\)。

\subsection{深度的存在性与有限性}
\label{b3:subsec:1601:02}

\begin{definition}\label{b3:def:depth}
设 \((R,\mathfrak m)\) 为交换 Noether 局部环，\(M\ne0\) 为有限生成的 \(R\) 模。\(\mathfrak m\) 中任一极大 \(M\) 正则序列的长度称为 \(M\) 在 \(R\) 上的\term[shendu]{深度}[depth]，记为 \(\operatorname{depth}_RM\)。约定 \(\operatorname{depth}_R0=+\infty\)；环的深度是其作为自身模的深度。
\end{definition}
\symidx{\(\operatorname{depth}_RM\)：局部模的深度}

定义的选择无关性由\cref{b3:thm:maximal-regular-length}保证，也说明非零有限生成模的深度是有限非负整数。固定极大理想的有限生成组后，它不超过生成元数。零模取无穷大的约定只用于公式，不把零模解释为可以在本章定义下具有任意长正则序列的对象。

\begin{corollary}\label{b3:cor:depth-zero-socle}
设 \((R,\mathfrak m)\) 为交换 Noether 局部环，\(k=R/\mathfrak m\) 为其剩余域，\(M\ne0\) 为有限生成的 \(R\) 模。以下条件等价：
\begin{equivalences}
\item \(\operatorname{depth}_RM=0\)；
\item \(\mathfrak m\in\operatorname{Ass}_RM\)；
\item \(\operatorname{Hom}_R(k,M)\ne0\)；
\item \(M\) 含有非零有限长度子模。
\end{equivalences}
\end{corollary}
\begin{proof}
前两个条件的等价已在\cref{b3:thm:maximal-regular-length}的证明中得到。一个同态 \(k\rightarrow M\) 由 \(1+\mathfrak m\) 的像决定，该像恰为被 \(\mathfrak m\) 零化的元素；非零时其零化子等于 \(\mathfrak m\)，故第二、第三个条件等价。这样的非零同态从简单模 \(k\) 出发，必为单射，其像是长度一的子模，因而得到第四个条件。

反过来，设 \(0\ne L\subseteq M\) 具有有限合成列 \(0=L_0\subsetneq L_1\subsetneq\cdots\subsetneq L_\ell=L\)。局部环上的每个简单商 \(L_j/L_{j-1}\) 都同构于 \(k\)，所以 \(\mathfrak mL_j\subseteq L_{j-1}\)。沿合成列逐步相乘得到 \(\mathfrak m^\ell L=0\)。取使 \(\mathfrak m^aL=0\) 的最小正整数 \(a\)，则 \(\mathfrak m^{a-1}L\ne0\)，其中的非零元素被 \(\mathfrak m\) 零化，于是给出非零同态 \(k\to M\)。这证明第四个条件蕴含第三个条件。
\end{proof}

这样的简单子模可能不唯一；把模中的所有简单子模合在一起，就得到下面的子模。

\begin{definition}\label{b3:def:socle}
设 \((R,\mathfrak m)\) 为交换局部环，\(k=R/\mathfrak m\) 为其剩余域，\(M\) 为 \(R\) 模。\(M\) 中全部简单子模之和称为 \(M\) 的\term[jizuo]{基座}[socle]，记为 \(\operatorname{Soc}_RM\)。在此局部环情形，有
\[
\operatorname{Soc}_RM=(0:_M\mathfrak m)
=\{\,u\in M\mid au=0\text{ 对所有 }a\in\mathfrak m\,\}.
\]
\end{definition}
\symidx{\(\operatorname{Soc}_RM\)：模的基座}

每个简单 \(R\) 模都同构于 \(k\)，因此被 \(\mathfrak m\) 零化；反过来，若 \(u\ne0\) 且 \(\mathfrak mu=0\)，则 \(Ru\cong k\) 是简单子模，证明了定义中两种描述相同。求值映射
\[
\operatorname{Hom}_R(k,M)\longrightarrow\operatorname{Soc}_RM,
\qquad f\longmapsto f(1+\mathfrak m)
\]
是 \(k\) 线性同构，其逆把 \(u\) 送到同态 \(a+\mathfrak m\mapsto au\)。所以当 \(R\) 还是 Noether 环、\(M\ne0\) 且有限生成时，\cref{b3:cor:depth-zero-socle}给出
\[
\operatorname{depth}_RM=0\quad\Longleftrightarrow\quad
\operatorname{Soc}_RM\ne0.
\]

例如，对域 \(k\)，章首的环 \(A=k[X,Y]_{(X,Y)}/(X^2,XY)\) 中，非零的 \(\bar X\) 被 \((\bar X,\bar Y)\) 零化，所以 \(\bar X\in\operatorname{Soc}_AA\)。它生成一个同构于剩余域 \(k\) 的长度一子模。这正是该环虽有一维支集，却找不到第一个正则元素的原因。

\begin{corollary}\label{b3:cor:regular-quotient-depth}
设 \((R,\mathfrak m)\) 为交换 Noether 局部环，\(M\ne0\) 为有限生成的 \(R\) 模，\(x\in\mathfrak m\) 在 \(M\) 上为非零因子。则
\[
\operatorname{depth}_R(M/xM)
=\operatorname{depth}_{R/(x)}(M/xM)
=\operatorname{depth}_RM-1.
\]
\end{corollary}
\begin{proof}
在 \(M/xM\) 上取极大正则序列，提升其元素到 \(\mathfrak m\)，并把 \(x\) 放在最前面，得到 \(M\) 上的极大正则序列，长度增加一。\(R\) 与 \(R/(x)\) 的作用在此商模上相同，两者的极大理想内元素及逐次商一一对应，故两种底环下深度相等。
\end{proof}

例如 \(k[X,Y]_{(X,Y)}\) 上的环本身深度为二，而模 \(k\) 深度为零。若 \(R\) 为任意非零 Artin 局部环，则非零有限生成模的基座非零，所以深度为零。这并不表示模很小：基座只保证正则序列无法迈出第一步。

\subsection{深度与维数的不等式}
\label{b3:subsec:1601:03}

\begin{proposition}\label{b3:prop:regular-quotient-dimension}
设 \((R,\mathfrak m)\) 为交换 Noether 局部环，\(M\ne0\) 为有限生成的 \(R\) 模，\(x\in\mathfrak m\) 在 \(M\) 上为非零因子。则
\[
\dim_R(M/xM)=\dim_RM-1.
\]
\end{proposition}
\begin{proof}
令 \(A=R/\operatorname{Ann}_RM\)，其极大理想仍记为 \(\mathfrak m\)。在任一素理想处，若 \(x\) 不在该素理想中，商模的局部化为零；若 \(x\) 在其中且 \(M\) 的局部化非零，Nakayama 引理说明商模的局部化非零。因此
\[
\operatorname{Supp}(M/xM)=\operatorname{Supp}M\cap V(x),\qquad
\dim(M/xM)=\dim A/(x).
\]
要证明维数恰好下降一，须分别排除下降不足一维和超过一维。由\cref{b3:prop:minimal-primes-associated}，\(\operatorname{Supp}M\) 的极小素理想都是关联素理想。\(x\) 为非零因子，故避开这些素理想：方程 \(x=0\) 不会保留支集的任何一个完整不可约分量。每条包含 \(x\) 的支集素链都可以在底部再补上一个不含 \(x\) 的极小素理想，这个包含严格，所以 \(\dim A/(x)\le\dim A-1\)。这保证确实至少下降一维。

还须证明一个方程至多使维数下降一。设 \(d'=\dim A/(x)\)，由\cref{b3:thm:system-of-parameters}在这个非零 Noether 局部商环选 \(d'\) 个元素，其生成理想的根为极大理想。提升到 \(A\)，再加上 \(x\)，得到由 \(d'+1\) 个元素生成的 \(\mathfrak m\) 准素理想。\(\mathfrak m\) 是该理想上唯一的素理想，故由\cref{b3:thm:height-theorem}有 \(\dim A=\operatorname{ht}\mathfrak m\le d'+1\)。这给出反向不等式 \(\dim A/(x)\ge\dim A-1\)，两式合并得等式。
\end{proof}

\begin{corollary}\label{b3:cor:depth-dimension-bound}
设 \((R,\mathfrak m)\) 为交换 Noether 局部环，\(M\ne0\) 为有限生成的 \(R\) 模。则
\[
0\le\operatorname{depth}_RM\le\dim_RM\le\dim R.
\]
\end{corollary}
\begin{proof}
沿一个长度为 \(t\) 的极大正则序列逐次取商，维数每次下降一，末端模仍非零，其维数至少为零。因此 \(t\le\dim_RM\)。最后一式来自支集是素谱的子集。
\end{proof}

深度可能严格小于维数。前面含有基座元素 \(\bar X\) 的环 \(A\) 维数为一、深度为零；维数仍看见 \(Y\) 方向的素理想链，正则性却已被基座中的非零元素阻止。一般地，沿极大正则序列每次取商，维数和深度都下降一，所以末端商的维数是 \(\dim_RM-\operatorname{depth}_RM\)。这个数表示正则元素已经无法再选取时，支集中还剩多少维。

上述上界是 \(\dim_RM\)，不能换成 \(\operatorname{depth}R\)。例如 \(A\) 自身深度为零，但模 \(N=A/(\bar X)\cong k[Y]_{(Y)}\) 中乘 \(\bar Y\) 单射，最终商为 \(k\)，所以 \(\operatorname{depth}_AN\ge1\)；又 \(\dim_AN=1\)，故其深度恰为一，大于环 \(A\) 的深度。

\begin{proposition}\label{b3:prop:nonreduced-positive-depth}
设 \(k\) 为域，\(a\ge2\) 为整数，令
\[
A=k[X,Y]_{(X,Y)}/(X^a),
\]
并以 \(\bar X,\bar Y\) 表示 \(X,Y\) 在 \(A\) 中的像。则 \(\bar Y\) 是 \(A\) 的正则元素，\(\dim A=\operatorname{depth}A=1\)，但 \(\bar X\ne0\) 且 \(\bar X^a=0\)，所以 \(A\) 不约化。
\end{proposition}
\begin{proof}
环 \(k[X,Y]/(X^a)\) 作为 \(k[Y]\) 模自由，基为 \(1,\bar X,\ldots,\bar X^{a-1}\)，所以乘 \(Y\) 单射。局部化保持单射，故乘 \(\bar Y\) 在 \(A\) 上仍单射。其商为
\[
A/(\bar Y)\cong k[X]/(X^a)\ne0,
\]
右边已是局部环，因为 \(X\) 幂零，唯一极大理想是 \((\bar X)\)。因此 \(\bar Y\) 正则，深度至少为一。上述商中 \(\bar X\) 非零，因为 \(a\ge2\)；于是它在 \(A\) 中也非零，而 \(\bar X^a=0\)。又 \((\bar X)\) 是幂零理想，且 \(A/(\bar X)\cong k[Y]_{(Y)}\) 约化，故 \(\sqrt{(0)}=(\bar X)\)。由\cref{b3:prop:dimension-basic-comparison}，\(\dim A=\dim k[Y]_{(Y)}=1\)。深度与维数的不等式再给出 \(\operatorname{depth}A=1\)。
\end{proof}

% !TeX root = ../../Book3.tex
% 纲要标题：### 16.2 Ext 刻画
% 纲要位置：计划纲要.md，第 2466 行。

\section{Ext 刻画}
\label{b3:sec:1602}

一个正则元素在模上作用为单射，却把剩余域零化。把这两种作用放进 Hom 的消解计算，短正合列便使首次非零的上同调次数改变一。正则序列的长度因此有另一种表达：从剩余域到模的 Ext 首次非零次数。局部化时，剩余域随素理想改变，还需要把这种计算与支集维数联系起来。


% 纲要内容（仅作写作计划，尚非教材正文）：
%
% * Ext 的两变量长正合列、零化作用、有限性与局部化
% * 以剩余域的 Ext 检测深度
% * 与正则序列定义的等价
% * Ischebeck 消失、关联素理想维数界与深度的局部化估计
% * 正则序列、Koszul 同调与 Ext 的深度刻画对照
%

\subsection{剩余域的 Ext 与深度检测}
\label{b3:subsec:1602:01}

\begin{definition}\label{b3:def:ext-minimum-background}
设 \(R\) 为交换环，\(N,M\) 为 \(R\) 模，选自由消解 \(P_\bullet\rightarrow N\rightarrow0\)，其微分记为 \(\mathrm{d}_i:P_i\rightarrow P_{i-1}\)（\(i\ge1\)）。取上链复形
\[
C^i=\operatorname{Hom}_R(P_i,M)\quad(i\ge0),\qquad
\delta^i(f)=f\circ \mathrm{d}_{i+1}.
\]
令 \(C^i=0\) 对 \(i<0\) 成立，\(\delta^{-1}=0\)。因 \(\mathrm{d}_i\mathrm{d}_{i+1}=0\)，有 \(\delta^{i+1}\delta^i=0\)。定义
\[
\operatorname{Ext}_R^i(N,M)=H^i(C^\bullet)
=\ker\delta^i/\operatorname{im}\delta^{i-1}\quad(i\ge0),
\]
并令负次数的 Ext 为零。
\end{definition}
\symidx{\(\operatorname{Ext}_R^i(N,M)\)：Ext 模}

由\cref{b3:lem:free-resolution-existence}，每个模都有自由消解。\cref{b3:lem:free-resolution-comparison}给出不同消解之间的比较映射及其链同伦唯一性；施加 \(\operatorname{Hom}_R(-,M)\) 后，同伦仍给出相同的上同调映射。因此 Ext 不依赖选择，且次数零自然同构于 \(\operatorname{Hom}_R(N,M)\)。同态 \(N'\rightarrow N\) 提升为消解的链映射后，经 Hom 反向诱导 Ext 映射；同态 \(M\rightarrow M'\) 则经后复合正向诱导映射。比较引理使这些映射与选择无关，并保持恒等映射及复合，所以 Ext 在第一变量反变、第二变量协变。

\begin{lemma}[Ext 的长正合列与零化作用]\label{b3:lem:ext-basic-properties}
设 \(R\) 为交换环，\(N,M\) 为 \(R\) 模。
\begin{enumerate}
\item 短正合列在 Ext 的第二变量给出协变长正合列，在第一变量给出反变长正合列。
\item 若 \(a\in R\) 满足 \(aN=0\) 或 \(aM=0\)，则 \(a\operatorname{Ext}_R^i(N,M)=0\) 对所有 \(i\ge0\) 成立。
\end{enumerate}
\end{lemma}
\begin{proof}
第二变量的短正合列经每个 \(\operatorname{Hom}_R(P_i,-)\) 仍正合，因为自由模的映射可对基逐个提升。于是得到短正合的上链复形。其连接映射这样构造：提升商复形中的一个圈，取其微分，结果落在子复形中且为圈；改变提升或原圈代表只改变一个边界。这也给出正合性：连接类为零恰好表示可以修正提升使其成为圈，其他两处分别对应圈来自子复形和边界能够提升。

第一变量的短正合列记为 \(0\rightarrow N'\rightarrow N\rightarrow N''\rightarrow0\)，给两端选自由消解 \(P',P''\)。令 \(Z'_{-1}=N'\)、\(Z_{-1}=N\)、\(Z''_{-1}=N''\)。递归构造 \(P\) 时，假设已有
\[
0\longrightarrow Z'_{i-1}\xrightarrow{j}Z_{i-1}
\xrightarrow{q}Z''_{i-1}\longrightarrow0.
\]
两端消解给出满射 \(\rho'_i:P'_i\rightarrow Z'_{i-1}\)、\(\rho''_i:P''_i\rightarrow Z''_{i-1}\)，记其核为 \(Z'_i,Z''_i\)。因 \(P''_i\) 自由，可取提升 \(\ell_i:P''_i\rightarrow Z_{i-1}\)，使 \(q\ell_i=\rho''_i\)。定义
\[
P_i=P'_i\oplus P''_i,\qquad
\rho_i(u,v)=j\rho'_i(u)+\ell_i(v).
\]
它为满射：先从 \(P''_i\) 提升 \(qz\)，余下的差属于 \(jZ'_{i-1}\)，再从 \(P'_i\) 提升该差。令 \(Z_i=\ker\rho_i\)。包含第一分量与投影第二分量使
\[
0\longrightarrow Z'_i\longrightarrow Z_i\longrightarrow Z''_i\longrightarrow0
\]
正合。当前这一层的构造放在三列图中如下：中排是刚选出的自由模，向下满射到上一层的三个核，向上则是本层新得到的三个核。各行均短正合，中排还分裂。
\[
\begin{tikzcd}[column sep=2em,row sep=2.5em]
0\arrow[r]&Z'_i\arrow[r,hook]\arrow[d,hook]
 &Z_i\arrow[r,two heads]\arrow[d,hook]
 &Z''_i\arrow[r]\arrow[d,hook]&0\\
0\arrow[r]&P'_i\arrow[r,hook]\arrow[d,two heads,"\rho'_i"']
 &P'_i\oplus P''_i\arrow[r,two heads]\arrow[d,two heads,"\rho_i"]
 &P''_i\arrow[r]\arrow[d,two heads,"\rho''_i"]&0\\
0\arrow[r]&Z'_{i-1}\arrow[r,hook,"j"']
 &Z_{i-1}\arrow[r,two heads,"q"']
 &Z''_{i-1}\arrow[r]&0
\end{tikzcd}
\]
构造 \(\rho_i\) 时，先用 \(\ell_i\) 提升右端，左端再修正误差。尤其，对 \(v\in Z''_i\)，\(\ell_i(v)\) 落在 \(jZ'_{i-1}\) 中，可取 \(u\) 使 \(j\rho'_i(u)=-\ell_i(v)\)，故 \((u,v)\in Z_i\) 提升 \(v\)；投影的核恰是 \(Z'_i\)。上排因此具有与下排相同的形式，可以作为下一层的输入。继续递归，取 \(\rho_0\) 为增广映射，并在 \(i\ge1\) 时把 \(\rho_i\) 与 \(Z_{i-1}\hookrightarrow P_{i-1}\) 复合，得到中间的自由消解 \(P\)。其包含与投影保持微分，组成逐项分裂的短正合消解
\[
0\longrightarrow P'\longrightarrow P\longrightarrow P''\longrightarrow0,
\qquad P_i=P'_i\oplus P''_i.
\]
于是每个次数有分裂短正合列
\[
0\longrightarrow\operatorname{Hom}_R(P''_i,M)
\longrightarrow\operatorname{Hom}_R(P_i,M)
\longrightarrow\operatorname{Hom}_R(P'_i,M)\longrightarrow0.
\]
前段圈提升的构造就给出第一变量的反变长正合列。这一逐项分裂的消解构造称为马蹄构造；连接映射与比较映射相容，因而所得长正合列自然。

若 \(aN=0\)，则 \(a\operatorname{id}_{P_\bullet}\) 与零映射都提升 \(N\) 上的零同态。\cref{b3:lem:free-resolution-comparison}给出链同伦 \(a\operatorname{id}=\mathrm{d}h+h\mathrm{d}\)。施加 Hom 后，乘法 \(a\) 在上同调上为零。若 \(aM=0\)，Hom 复形各项上的乘法 \(a\) 已经为零。
\end{proof}

\begin{proposition}[有限生成的 Ext 与局部化]\label{b3:prop:finite-ext-localization}
设 \(R\) 为交换 Noether 环，\(N,M\) 为有限生成的 \(R\) 模，\(i\ge0\) 为整数。则 \(\operatorname{Ext}_R^i(N,M)\) 是有限生成的 \(R\) 模，并且任意乘法集 \(S\subseteq R\) 都给出自然同构
\[
S^{-1}\operatorname{Ext}_R^i(N,M)
\cong\operatorname{Ext}_{S^{-1}R}^i(S^{-1}N,S^{-1}M).
\]
\end{proposition}
\begin{proof}
因 \(N\) 有限生成，可先选有限秩自由模到 \(N\) 的满射。每个核都由 Noether 性有限生成，所以可以逐次选有限秩自由模满射，得到各项均为有限秩的自由消解 \(P\)。每个 \(\operatorname{Hom}_R(P_i,M)\) 是有限生成模，其核和像也有限生成，故上同调有限生成。对有限秩自由模有
\[
S^{-1}\operatorname{Hom}_R(P_i,M)
\cong\operatorname{Hom}_{S^{-1}R}(S^{-1}P_i,S^{-1}M),
\]
且这些同构保持微分。局部化又保持消解的正合性、核及像，因此取上同调得到所述自然同构。
\end{proof}

\subsection{Ext 刻画与正则序列定义的等价性}
\label{b3:subsec:1602:02}

\begin{theorem}\label{b3:thm:depth-ext}
设 \((R,\mathfrak m)\) 为交换 Noether 局部环，\(k=R/\mathfrak m\) 为其剩余域，\(M\ne0\) 为有限生成的 \(R\) 模。则
\[
\operatorname{depth}_RM
=\inf\{\,i\ge0\mid\operatorname{Ext}_R^i(k,M)\ne0\,\}.
\]
特别地，右侧集合非空，其最小值等于正则序列定义的深度。
\end{theorem}
\begin{proof}
对 \(t=\operatorname{depth}_RM\) 归纳。\(t=0\) 时，由\cref{b3:cor:depth-zero-socle}，\(\operatorname{Ext}_R^0(k,M)=\operatorname{Hom}_R(k,M)\ne0\)。设 \(t>0\)，选 \(x\in\mathfrak m\) 为 \(M\) 非零因子，令 \(\bar M=M/xM\)。由\cref{b3:cor:regular-quotient-depth}，\(\operatorname{depth}_R\bar M=t-1\)。

短正合列 \(0\rightarrow M\xrightarrow{x}M\rightarrow\bar M\rightarrow0\) 给出 Ext 长正合列。因为 \(xk=0\)，\cref{b3:lem:ext-basic-properties}说明所有 Ext 上的乘法 \(x\) 为零，故每个 \(i\ge0\) 有
\[
0\longrightarrow\operatorname{Ext}_R^i(k,M)
\longrightarrow\operatorname{Ext}_R^i(k,\bar M)
\longrightarrow\operatorname{Ext}_R^{i+1}(k,M)\longrightarrow0.
\]
又因 \(x\) 在 \(M\) 上单射，\(\operatorname{Hom}_R(k,M)=0\)。归纳假设说中间项在 \(i<t-1\) 为零、在 \(i=t-1\) 非零。上述短正合列遂给出 \(\operatorname{Ext}_R^i(k,M)=0\) 对 \(i<t\) 成立，并给出
\[
\operatorname{Ext}_R^t(k,M)\cong\operatorname{Ext}_R^{t-1}(k,\bar M)\ne0.
\]
这完成归纳。
\end{proof}

例如在 \(R=k[X,Y]_{(X,Y)}\) 上，用 \(X,Y\) 的 Koszul 消解计算 \(\operatorname{Ext}_R^i(k,R)\)。对偶复形的两个矩阵是 \(\binom XY\) 与 \((-Y\ \ X)\)，其前两处上同调为零，最后的余核为 \(k\)。所以深度为二。在 \(R=k[\varepsilon]/(\varepsilon^2)\) 上，映射 \(k\rightarrow R\)、\(1\mapsto\varepsilon\) 已在次数零非零，故深度为零，不必先构造整个无限消解。

\ProseParagraphBreak
在第一个例子中，若交换两个变量，\(R\) 自身已是自由模，可以用只含零次项的消解计算，故 \(\operatorname{Ext}_R^{i>0}(R,k)=0\)。这与 \(\operatorname{Ext}_R^2(k,R)\cong k\) 同时成立：深度判据把 \(k\) 放在第一变量、被检测的模放在第二变量。后面的投射维数判据检测的是被消解的第一变量。Tor 的变量对称性不能照搬给 Ext。

\subsection{Ischebeck 消失与深度的局部化估计}
\label{b3:subsec:1602:03}

设 \((R,\mathfrak m)\) 为交换 Noether 局部环，\(k=R/\mathfrak m\)。局部化 Ext 的公式不能把闭点的剩余域直接变成其他素理想处的剩余域。若 \(\mathfrak p\subsetneq\mathfrak m\)，则 \(k_{\mathfrak p}=0\)，而 \(\kappa(\mathfrak p)\) 非零。比较两个地方的深度，需要额外的维数估计。Ischebeck 的消失定理把第一变量的支集维数与第二变量的深度联系起来。\footnote{弗里德里希·格哈特·伊舍贝克（Friedrich Gerhart Ischebeck，1940--2022），德国数学家，主要研究交换代数和同调代数，著有《Einladung zur Zahlentheorie》。}

\begin{lemma}\label{b3:lem:ischebeck-vanishing}
设 \((R,\mathfrak m)\) 为交换 Noether 局部环，\(M,N\ne0\) 为有限生成的 \(R\) 模，\(t=\operatorname{depth}_RM\)、\(s=\dim_RN\)。则
\[
\operatorname{Ext}_R^i(N,M)=0\qquad(0\le i<t-s).
\]
\end{lemma}
\begin{proof}
对 \(s\) 归纳，先说明如何归约到素循环模。任一非零有限生成模都有有限素因子滤过：在非零商中选一个关联素理想，便得到同构于 \(R/\mathfrak q\) 的子模，再在余下的商中重复；相应的子模严格升链由 Noether 性终止。将所得滤过记为
\[
0=N_\ell\subset N_{\ell-1}\subset\cdots\subset N_0=N,
\qquad N_j/N_{j+1}\cong R/\mathfrak q_j.
\]
各因子的支集包含于 \(\operatorname{Supp}N\)，所以 \(\dim(R/\mathfrak q_j)\le s\)。若对某个次数 \(i\) 所有因子的 Ext 都为零，从 \(N_\ell=0\) 开始对滤过长度归纳。短正合列
\(0\rightarrow N_{j+1}\rightarrow N_j\rightarrow R/\mathfrak q_j\rightarrow0\)
的第一变量长正合列给出
\[
\operatorname{Ext}_R^i(R/\mathfrak q_j,M)
\longrightarrow\operatorname{Ext}_R^i(N_j,M)
\longrightarrow\operatorname{Ext}_R^i(N_{j+1},M).
\]
外两项为零便使中项为零，逐层推出 \(\operatorname{Ext}_R^i(N,M)=0\)。因此，只需对 \(N=R/\mathfrak q\) 证明指定范围内的消失。

当 \(s=0\) 时，各因子的素理想只能为 \(\mathfrak m\)，因子就是 \(k\)。\cref{b3:thm:depth-ext}给出各因子的 Ext 在 \(i<t\) 时为零，前段长度归纳完成基例。设 \(s>0\)。若 \(\dim(R/\mathfrak q)<s\)，由维数归纳假设得到消失；其消失范围至少包含 \(i<t-s\)。剩下 \(\dim(R/\mathfrak q)=s\) 的情形。选 \(x\in\mathfrak m\setminus\mathfrak q\)，则乘法 \(x\) 在 \(R/\mathfrak q\) 上单射，有
\[
0\longrightarrow R/\mathfrak q\xrightarrow{x}R/\mathfrak q
\longrightarrow R/(\mathfrak q,x)\longrightarrow0.
\]
由\cref{b3:prop:regular-quotient-dimension}，末项维数为 \(s-1\)。归纳假设使其 Ext 在次数小于 \(t-s+1\) 时为零。对 \(0\le i<t-s\)，第一变量长正合列中有
\[
\operatorname{Ext}_R^i(R/\mathfrak q,M)
\xrightarrow{x}\operatorname{Ext}_R^i(R/\mathfrak q,M)
\longrightarrow\operatorname{Ext}_R^{i+1}(R/(\mathfrak q,x),M)=0.
\]
因此乘法 \(x\) 为满射。\cref{b3:prop:finite-ext-localization}说明该 Ext 模有限生成；因 \(x\in\mathfrak m\)，Nakayama 引理使其为零。再按素因子滤过的长度归纳，便得到原模 \(N\) 的结论。
\end{proof}

\begin{corollary}\label{b3:cor:associated-depth-bound}
设 \((R,\mathfrak m)\) 为交换 Noether 局部环，\(M\ne0\) 为有限生成的 \(R\) 模。若 \(\mathfrak q\in\operatorname{Ass}_RM\)，则
\[
\operatorname{depth}_RM\le\dim(R/\mathfrak q).
\]
\end{corollary}
\begin{proof}
关联素理想给出非零单射 \(R/\mathfrak q\rightarrow M\)。若维数小于深度，上述引理将使 \(\operatorname{Hom}_R(R/\mathfrak q,M)=0\)，与该单射矛盾。
\end{proof}

\begin{theorem}\label{b3:thm:depth-localization}
设 \((R,\mathfrak m)\) 为交换 Noether 局部环，\(M\ne0\) 为有限生成的 \(R\) 模，\(\mathfrak p\in\operatorname{Supp}_RM\)。则
\[
\operatorname{depth}_RM
\le\operatorname{depth}_{R_{\mathfrak p}}M_{\mathfrak p}
+\dim(R/\mathfrak p).
\]
\end{theorem}
\begin{proof}
在 \(\mathfrak p\) 中连续选取 \(M\) 正则元素，直至不能延长，记所得序列长度为 \(r\)，末端商为 \(L\)。每加进一个正则元素，子模 \((x_1,\ldots,x_j)M\) 严格增大，所以 \(M\) 的 Noether 性保证过程终止。\(\mathfrak p\) 中所有元素都是 \(L\) 的零因子；而 \(\operatorname{Ass}_RL\) 有限，故由有限素理想回避，存在 \(\mathfrak q\in\operatorname{Ass}_RL\)，使 \(\mathfrak p\subseteq\mathfrak q\)。\cref{b3:cor:regular-quotient-depth}和\cref{b3:cor:associated-depth-bound}给出
\[
\operatorname{depth}_RM=r+\operatorname{depth}_RL
\le r+\dim(R/\mathfrak q)\le r+\dim(R/\mathfrak p).
\]
原序列局部化后仍逐次单射，而且所有元素都在 \(\mathfrak pR_{\mathfrak p}\) 中。\(M_{\mathfrak p}\ne0\) 且有限生成，Nakayama 引理保证每个商非零，所以这是一个长度 \(r\) 的 \(M_{\mathfrak p}\) 正则序列。于是 \(r\le\operatorname{depth}_{R_{\mathfrak p}}M_{\mathfrak p}\)，代入便得结论。
\end{proof}

有限生成条件使关联素理想集及 Ext 模受 Noether 性控制，也使 Nakayama 引理适用。例如，设 \(R\) 为离散赋值环，\(\pi\) 为其统一化元，\(K\) 为分式域。乘法 \(\pi\) 在 \(K\) 上既单射又满射，因而 \(K/\pi K=0\)；它不满足正则序列的末端非零条件。

\ProseParagraphBreak
\cref{b3:thm:maximal-regular-length,b3:thm:depth-ext}将正则序列的长度与两种同调计算联系起来，\cref{b3:cor:depth-zero-socle}给出深度零的等价条件。这些刻画集中列于\cref{b3:tab:depth-characterizations}。

\begin{center}
\begin{minipage}{\textwidth}
\makeatletter
\def\@captype{table}
\makeatother
\centering
\caption[深度的三种刻画]{深度的三种刻画。设 \((R,\mathfrak m)\) 为交换 Noether 局部环，\(M\ne0\) 为有限生成的 \(R\) 模，\(k=R/\mathfrak m\)，\(y=(y_1,\ldots,y_e)\) 为 \(\mathfrak m\) 的任意有限生成组。前三行给出同一个数值，末行列出该数值为零时的等价条件。}
\label{b3:tab:depth-characterizations}
\begin{tabularx}{\textwidth}{@{}>{\raggedright\arraybackslash}p{0.25\textwidth}>{\raggedright\arraybackslash}X@{}}
\toprule
刻画 & 数值或条件\\
\midrule
正则序列 & \(\mathfrak m\) 中任一极大 \(M\) 正则序列的长度\\
Koszul 同调 & \(e-\max\{\,i\mid H_i(y;M)\ne0\,\}\)\\
Ext & \(\min\{\,i\ge0\mid\operatorname{Ext}_R^i(k,M)\ne0\,\}\)\\
深度零 & \(\mathfrak m\in\operatorname{Ass}_RM\ \Longleftrightarrow\ \operatorname{Soc}_RM\ne0\)\\
\bottomrule
\end{tabularx}
\end{minipage}
\end{center}

% !TeX root = ../../Book3.tex
% 纲要标题：### 16.3 深度引理与合冲模
% 纲要位置：工作记录/计划纲要.md，第 2501 行。
% 纲要细目（写作范围）：
% * 短正合列中的深度不等式与零模的无穷大约定
% * 合冲模的深度下界及其不能充当上界的例子
% * 用消解逐步比较深度

\section{深度引理与合冲模}
\label{b3:sec:1603}

短正合列把三个模联系起来，深度却通常只满足不等式。首次非零的 Ext 及其连接同态可以说明何时达到等号，也能追踪商去一个非零因子后深度的变化。


% 纲要内容（仅作写作计划，尚非教材正文）：
%
% * 短正合列中的深度不等式与零模的无穷大约定
% * 合冲模的深度下界及其不能充当上界的例子
% * 用消解逐步比较深度
%

\subsection{短正合列中的深度不等式}
\label{b3:subsec:1603:01}

\begin{theorem}[深度引理]\label{b3:thm:depth-lemma}
设 \((R,\mathfrak m)\) 为交换 Noether 局部环，\(k=R/\mathfrak m\) 为其剩余域，
\[
0\longrightarrow A\longrightarrow B\longrightarrow C\longrightarrow0
\]
为有限生成的 \(R\) 模的短正合列。约定零模深度为 \(+\infty\)，并采用
\[
(+\infty)+1=(+\infty)-1=+\infty,
\qquad \min\{+\infty,t\}=t
\quad(t\in\mathbb Z\cup\{+\infty\}).
\]
则
\begin{align*}
\operatorname{depth}B&\ge\min\{\operatorname{depth}A,\operatorname{depth}C\},\\
\operatorname{depth}A&\ge\min\{\operatorname{depth}B,\operatorname{depth}C+1\},\\
\operatorname{depth}C&\ge\min\{\operatorname{depth}A-1,\operatorname{depth}B\}.
\end{align*}
此外，若 \(\operatorname{depth}B>\operatorname{depth}C\) 且 \(C\ne0\)，则 \(\operatorname{depth}A=\operatorname{depth}C+1\)；若 \(\operatorname{depth}A<\operatorname{depth}B\) 且 \(A\ne0\)，则 \(\operatorname{depth}C=\operatorname{depth}A-1\)。
\end{theorem}
\begin{proof}
令 \(E^i(X)=\operatorname{Ext}_R^i(k,X)\)，并对 \(i<0\) 置 \(E^i(X)=0\)。短正合列给出
\[
\cdots\longrightarrow E^{i-1}(C)\longrightarrow E^i(A)
\longrightarrow E^i(B)\longrightarrow E^i(C)
\longrightarrow E^{i+1}(A)\longrightarrow\cdots.
\]
由\cref{b3:thm:depth-ext}，深度就是首次非零次数。三条不等式分别读这条长列中以 \(E^i(B)\)、\(E^i(A)\)、\(E^i(C)\) 为中项的三个片段。若 \(i\) 小于 \(A,C\) 的两个深度，\(E^i(A)=E^i(C)=0\)，故 \(E^i(B)=0\)，得到第一式。若 \(i\) 小于 \(\operatorname{depth}B\) 与 \(\operatorname{depth}C+1\)，则 \(E^{i-1}(C)=E^i(B)=0\)，得到第二式；同样由 \(E^i(B)=E^{i+1}(A)=0\) 得第三式。

若 \(c=\operatorname{depth}C<\operatorname{depth}B\)，则非零模 \(E^c(C)\) 单射到 \(E^{c+1}(A)\)。第二式又给出 \(\operatorname{depth}A\ge c+1\)，故恰等于 \(c+1\)。若 \(a=\operatorname{depth}A<\operatorname{depth}B\)，则长列使 \(E^{a-1}(C)\) 满射到非零的 \(E^a(A)\)；这也表明 \(a\ne0\)。第三式给出反向界，故深度恰为 \(a-1\)。这两个等式都来自连接映射：中项尚未出现非零 Ext 时，一端的首次非零类迫使另一端在相邻次数非零。零模的情形可从同构直接核对。
\end{proof}

三个下界都不能普遍改成等式。例如分裂列 \(0\rightarrow A\rightarrow A\oplus C\rightarrow C\rightarrow0\) 中，\(B\) 的深度确为两端较小者；但设 \(k\) 为域，对 \(R=k[X,Y]_{(X,Y)}\) 的列 \(0\rightarrow R\xrightarrow{X}R\rightarrow R/(X)\rightarrow0\)，两端深度分别为二和一，中项深度为二。

\subsection{合冲模的深度}
\label{b3:subsec:1603:02}

\begin{definition}\label{b3:def:syzygy-local}
设 \(R\) 为交换 Noether 环，\(M\) 为有限生成的 \(R\) 模，给定有限秩自由模 \(F\) 到 \(M\) 的满射。其核 \(\Omega M=\ker(F\rightarrow M)\) 称为这次自由展示所得的\term[hechongmo]{合冲模}[syzygy module]。重复对核选择自由满射，得到更高合冲模 \(\Omega^jM\)，其中 \(\Omega^0M=M\)。
\end{definition}

合冲模依赖所选展示。例如把 \(F\twoheadrightarrow M\) 换成 \(F\oplus R\twoheadrightarrow M\)，让新增分量映到零，核就从 \(\Omega M\) 变成 \(\Omega M\oplus R\)。当前记号因此总是指选定展示的核；本节的深度估计对每个选择都成立。下一章将用局部极小生成组固定更精细的消解。

\begin{corollary}\label{b3:cor:syzygy-depth}
设 \((R,\mathfrak m)\) 为交换 Noether 局部环，\(M\) 为有限生成的 \(R\) 模，\(j\ge0\) 为整数。在 \(M\) 的任一各项均为有限秩自由模的消解中，第 \(j\) 个合冲模满足
\[
\operatorname{depth}_R\Omega^jM
\ge\min\{\operatorname{depth}R,\operatorname{depth}_RM+j\}.
\]
若 \(M\ne0\)、\(\operatorname{depth}_RM<\operatorname{depth}R\)，则第一合冲模满足 \(\operatorname{depth}\Omega M=\operatorname{depth}M+1\)。
\end{corollary}
\begin{proof}
非零有限自由模 \(F=R^n\) 与 \(R\) 有相同深度，因为正则性在每个直和分量逐项检验。对 \(0\rightarrow\Omega M\rightarrow F\rightarrow M\rightarrow0\) 应用深度引理第二式，得到 \(\operatorname{depth}\Omega M\ge\min\{\operatorname{depth}R,\operatorname{depth}M+1\}\)；每换成一个更高的核，都可以重复这一步，逐次归纳即得下界。若两深度严格不等，深度引理的第一条附加等式给出精确值。遇到零合冲模时深度为无穷大，下界仍成立；这个约定不表示零模存在无限长的非退化正则序列。
\end{proof}

这个估计只是下界，\(\operatorname{depth}R\) 并不是合冲模实际深度的上界。设 \(k\) 为域。在 \(R=k[X,Y]_{(X,Y)}\) 上，剩余域 \(k\) 深度为零，其第一合冲模 \(\mathfrak m=(X,Y)\) 深度为一。第二合冲模由 \((-Y,X)\) 生成，与 \(R\) 同构，深度为二。这把前章 Koszul 消解的三个次数与深度的逐步增加对应起来。

\subsection{沿消解逐步比较深度}
\label{b3:subsec:1603:03}

当 \(\operatorname{depth}M+j\ge\operatorname{depth}R\) 时，上述下界不再随 \(j\) 增长；下面的模却会有高于环本身的深度。

\ProseParagraphBreak

例如，设 \(k\) 为域，取
\[
R=k[x,y]_{(x,y)}/(x^2,xy),\qquad M=R/(y).
\]
以下仍用 \(x,y\) 表示它们在 \(R\) 中的像。非零元 \(x\) 被极大理想 \((x,y)\) 零化，故 \(\operatorname{depth}R=0\)。自然满射 \(R\to M\) 的核为 \((y)\)。由于 \(xy=0\)，而乘 \(y\) 在 \(R/(x)\cong k[y]_{(y)}\) 上单射，\(\operatorname{Ann}_R(y)=(x)\)，所以
\[
\Omega M=(y)\cong R/(x)\cong k[y]_{(y)}.
\]
在该模上乘 \(y\) 单射，其支集维数为一，因而 \(\operatorname{depth}_R\Omega M=1>\operatorname{depth}R\)。消解长度与深度之间的精确关系还需要另外证明。

以最左端自由模的深度下界 \(\operatorname{depth}R\) 为起点，沿消解向右逐次使用深度引理，所得深度下界每步至多降低一。因此一个长度 \(n\) 的消解给出 \(\operatorname{depth}R-n\) 这个下界，也就对消解长度施加了必要限制。

\begin{proposition}\label{b3:prop:finite-resolution-depth-bound}
设 \((R,\mathfrak m)\) 为交换 Noether 局部环，\(M\ne0\) 为有限生成的 \(R\) 模，\(n\ge0\) 为整数。若有长度 \(n\)、各项均为有限秩自由模的消解
\[
0\longrightarrow F_n\longrightarrow\cdots\longrightarrow F_0
\longrightarrow M\longrightarrow0,
\]
则 \(\operatorname{depth}_RM\ge\operatorname{depth}R-n\)。
\end{proposition}
\begin{proof}
若 \(n=0\)，\(M\) 自由，结论成立。一般令 \(L=\ker(F_0\rightarrow M)\)。若 \(L=0\)，结论仍由自由性得出；否则尾部是 \(L\) 的长度 \(n-1\) 自由消解，归纳给出 \(\operatorname{depth}L\ge\operatorname{depth}R-(n-1)\)。深度引理第三式于是给出
\[
\operatorname{depth}M\ge
\min\{\operatorname{depth}L-1,\operatorname{depth}R\}
\ge\operatorname{depth}R-n.
\]
\end{proof}

例如，设 \(k\) 为域，取 \(R=k[X,Y,Z]_{(X,Y,Z)}\)、\(M=R/(X,Y)\)。Koszul 消解长为二，上式给出深度至少一；余下的 \(Z\) 确为正则元素，而商的维数为一，故深度恰为一。一般地，上式也写成 \(n\ge\operatorname{depth}R-\operatorname{depth}M\)，但这里的 \(n\) 是给定消解的长度，未必最短。对有限投射维数模，下一章将用极小消解证明 Auslander--Buchsbaum 公式，把这个必要界改进为投射维数的精确等式。

% !TeX root = ../../Book3.tex
% 纲要标题：### 16.4 Serre 条件
% 纲要位置：工作记录/计划纲要.md，第 2507 行。
% 纲要细目（写作范围）：
% * 模的 S_k 条件
% * 局部化后的深度要求
% * S_0、S_1、S_2 的层级与嵌入分量的联系
% * 一般 S_k 正则截面的下降；极大点数值不足与非约化反例
% ---

\section{Serre 条件}
\label{b3:sec:1604}

深度在一个点的值不能代替所有局部化处的检查。Serre 条件按各素点的局部维数提出统一下界，使闭点附近的退化与低余维处的性质能够分别比较。


% 纲要内容（仅作写作计划，尚非教材正文）：
%
% * 模的 S_k 条件
% * 局部化后的深度要求
% * S_0、S_1、S_2 的层级与嵌入分量的联系
%
% ---
%

\subsection{模的 \texorpdfstring{\(S_k\)}{Sk} 条件}
\label{b3:subsec:1604:01}

\begin{definition}\label{b3:def:serre-sk}
设 \(R\) 为交换 Noether 环，\(M\) 为有限生成的 \(R\) 模，\(k\ge0\) 为整数。若对每个 \(\mathfrak p\in\operatorname{Supp}_RM\) 都有
\[
\operatorname{depth}_{R_{\mathfrak p}}M_{\mathfrak p}
\ge\min\{k,\dim_{R_{\mathfrak p}}M_{\mathfrak p}\},
\]
则称 \(M\) 满足\term[serre-tiaojian]{Serre 条件}[Serre's condition] \((S_k)\)。对 \(M=R\) 采用同一定义，称环 \(R\) 满足 \((S_k)\)。零模的条件为空。
\end{definition}
\symidx{\((S_k)\)：Serre 深度条件}

深度不超过局部模的维数，所以在 \(\dim M_{\mathfrak p}\le k\) 的点，\((S_k)\) 要求深度恰等于该维数；在局部维数至少为 \(k\) 的点，则只要求深度至少为 \(k\)。维数为零的点没有额外要求。

这里采用局部模的维数，不能把它统一换成 \(\dim R_{\mathfrak p}\)：模的支集可能只有环境的一部分。例如，设 \(k\) 为域，\(R=k[X,Y]_{(X,Y)}\) 上的 \(M=R/(X)\) 只支撑在直线 \(X=0\) 上。在闭点，\(M\cong k[Y]_{(Y)}\) 的深度与模维数都为一；在这条线的一般点，局部模为 \(k(Y)\)，深度与维数都为零。因此它满足全部 Serre 条件，而不必具有深度二。

\subsection{局部化后的深度要求}
\label{b3:subsec:1604:02}

\begin{proposition}\label{b3:prop:serre-localization}
设 \(R\) 为交换 Noether 环，\(M\) 为有限生成的 \(R\) 模，\(k\ge0\) 为整数，\(S\subseteq R\) 为乘法集。
\begin{enumerate}
\item 若 \(M\) 满足 \((S_k)\)，则 \(S^{-1}M\) 在 \(S^{-1}R\) 上满足 \((S_k)\)。
\item \(M\) 满足 \((S_k)\) 当且仅当每个极大理想处的 \(M_{\mathfrak m}\) 满足 \((S_k)\)。
\item \((S_{k+1})\) 蕴含 \((S_k)\)，而每个有限生成模都满足 \((S_0)\)。
\end{enumerate}
\end{proposition}
\begin{proof}
局部化环的素理想对应于原环中避开乘法集的素理想，且在这些素理想处再局部化所得的环和模与直接局部化相同。因此深度与模维数都相同，第一条成立。任意素理想包含于某个极大理想，在该极大理想处先局部化再取原素理想，对象仍相同，得到第二条的反向；正向来自第一条。第三条直接比较最小值，局部环上非零有限生成模的深度非负。
\end{proof}

注意第二条要求极大理想处的局部模本身满足 \((S_k)\)，因而已经检查其全部素理想处的条件。只检查原环各极大理想处的一条深度数值，并不能自动完成定义中的所有检查。

\subsection{\texorpdfstring{\(S_1\)}{S1}、\texorpdfstring{\(S_2\)}{S2} 与嵌入分量}
\label{b3:subsec:1604:03}

\begin{theorem}\label{b3:thm:serre-s1-associated}
设 \(R\) 为交换 Noether 环，\(M\ne0\) 为有限生成的 \(R\) 模。则 \(M\) 满足 \((S_1)\) 当且仅当
\[
\operatorname{Ass}_RM=\operatorname{Min}(\operatorname{Supp}_RM),
\]
即其关联素理想全是支集的极小素理想。
\end{theorem}
\begin{proof}
对每个 \(\mathfrak p\in\operatorname{Supp}M\)，\(M_{\mathfrak p}\) 是非零有限生成的 \(R_{\mathfrak p}\) 模。由\cref{b3:cor:depth-zero-socle}，其深度为零当且仅当局部极大理想 \(\mathfrak pR_{\mathfrak p}\) 是关联素理想；再用\cref{b3:prop:associated-localization}及素理想的收缩对应，这又等价于 \(\mathfrak p\in\operatorname{Ass}_RM\)。另一方面，\(\dim M_{\mathfrak p}=0\) 当且仅当支集中没有严格小于 \(\mathfrak p\) 的素理想，即 \(\mathfrak p\) 是支集的极小元。

\((S_1)\) 只要求局部维数为正的各点深度不为零。因此它恰好排除非极小的关联素理想。由\cref{b3:prop:minimal-primes-associated}，支集的极小素理想总是关联素理想，所以排除这些嵌入关联素理想就得到所列集合等号。
\end{proof}

对域 \(k\)，第四章的两轴环 \(k[x,y]/(xy)\) 只有由 \((x)\)、\((y)\) 给出的两个极小关联素理想，所以满足 \((S_1)\)。而 \(k[x,y]/(x^2,xy)\) 中非零的 \(\bar x\) 被 \((x,y)\) 零化，其零化子恰为这个极大理想；它严格包含极小素理想 \((x)\)，所以该环不满足 \((S_1)\)。在这个闭点，局部维数仍为一，深度却已经降为零。

因此，对交换 Noether 环 \(R\) 上的有限生成模 \(M\)，这三层条件可以写成
\[
\begin{aligned}
(S_0)&:\quad\text{自动成立},\\
(S_1)&:\quad\text{没有嵌入关联素理想},\\
(S_2)&:\quad(S_1)\text{，并且在}\;\dim M_{\mathfrak p}\ge2
\;\text{的各点有}\;\operatorname{depth}M_{\mathfrak p}\ge2.
\end{aligned}
\]
它们满足 \((S_2)\Rightarrow(S_1)\Rightarrow(S_0)\)。这里 \((S_2)\) 仍包含一维局部模的深度至少为一的要求，不能只检查二维及更高维的点。

\ProseParagraphBreak

\((S_1)\) 与 \((S_2)\) 的差别可以在一个简单的理想上看见。设 \(k\) 为域，取 \(R=k[X,Y]_{(X,Y)}\)、\(M=(X,Y)\)。\(M\) 为整环中的非零理想，所有非零元的零化子都是零，所以它无挠，\(\operatorname{Ass}M=\{0\}\)，且支集为整个 \(\operatorname{Spec}R\)。它满足 \((S_1)\)；但前节算出其深度为一，模维数为二，故不满足 \((S_2)\)。排除嵌入关联素理想只保证正维局部点有第一层非零因子，并不保证二维及更高维处还能继续正则取商。

\begin{proposition}\label{b3:prop:s2-regular-section-s1}
设 \(R\) 为交换 Noether 环，\(s\ge1\) 为整数，\(M\) 为满足 \((S_s)\) 的有限生成的 \(R\) 模，\(x\in R\) 在 \(M\) 上为非零因子。则 \(M/xM\) 满足 \((S_{s-1})\)。特别地，满足 \((S_2)\) 的模沿非零因子取商后满足 \((S_1)\)。
\end{proposition}
\begin{proof}
局部化保持 \(x\) 的乘法单射。若 \(M_{\mathfrak p}\ne0\) 且 \(x\in\mathfrak p\)，Nakayama 引理保证 \(M_{\mathfrak p}/xM_{\mathfrak p}\ne0\)；若 \(x\notin\mathfrak p\)，则 \(x\) 成为单位，商模的局部化为零。因此
\[
\operatorname{Supp}(M/xM)=\operatorname{Supp}M\cap V(x).
\]
取这个支集中的 \(\mathfrak p\)，\(x/1\) 就是非零有限生成的 \(R_{\mathfrak p}\) 模 \(M_{\mathfrak p}\) 上的正则元素。记 \(d=\dim M_{\mathfrak p}\)、\(t=\operatorname{depth}M_{\mathfrak p}\)，由\cref{b3:prop:regular-quotient-dimension,b3:cor:regular-quotient-depth}，商的局部维数为 \(d-1\)，深度为 \(t-1\)。其中 \(d\ge1\)，而 \((S_s)\) 给出 \(t\ge\min\{s,d\}\)，故
\[
t-1\ge\min\{s,d\}-1=\min\{s-1,d-1\}.
\]
这正是商模在 \(\mathfrak p\) 处所需的 \((S_{s-1})\) 不等式。若商模为零，条件为空，结论也成立。
\end{proof}

因此 \((S_2)\) 不只排除当前模的嵌入部分，也控制正则截面不会新生这种部分。它仍不是正规性的完整条件；\chapref{b3:ch:20}将另外加入高度不超过一的各素点处的正则性，并证明两种要求如何共同刻画正规环。

\begin{proposition}\label{b3:prop:serre-condition-counterexamples}
非零交换 Artin 环满足所有 Serre 条件 \((S_s)\)，其中 \(s\ge0\) 为整数，但不必约化。例如，对任意域 \(k\)，环 \(k[\varepsilon]/(\varepsilon^2)\) 满足全部这些条件而不约化。

另设 \(k\) 为域，\(R=k[X,Y,Z]_{(X,Y,Z)}\)、\(\mathfrak m=(X,Y,Z)R\)、\(M=R\oplus R/(X,Y)\)。则 \(\operatorname{depth}_RM=1\)，但在素理想 \(\mathfrak p=(X,Y)R\) 处，有
\[
\operatorname{depth}_{R_{\mathfrak p}}M_{\mathfrak p}=0,
\qquad \dim_{R_{\mathfrak p}}M_{\mathfrak p}=2.
\]
因此 \(M\) 满足闭点自身的 \((S_1)\) 深度下界，却不满足 \((S_1)\)。
\end{proposition}
\begin{proof}
交换 Artin 环由\cref{b3:thm:artin-finite-length}是 Noether 环，且由\cref{b3:lem:artin-primes-radical}，每个素理想都极大。因此在每个素点，局部环的维数为零，Serre 条件只要求深度非负，自动成立。在 \(k[\varepsilon]/(\varepsilon^2)\) 中，\(\varepsilon\ne0\) 而 \(\varepsilon^2=0\)，所以约化性不能由这些深度条件推出。

在所列 \(R\) 与 \(M\) 中，\(R/(X,Y)\cong k[Z]_{(Z)}\)，所以 \(Z\) 在两个直和分量上都为非零因子，也就在 \(M\) 上为非零因子。商模
\[
M/ZM\cong R/(Z)\oplus k
\]
含有非零元素 \((0,1)\)，其零化子为 \(\mathfrak m\)。由\cref{b3:cor:depth-zero-socle}，该商模深度为零；再由\cref{b3:cor:regular-quotient-depth}，\(\operatorname{depth}_RM=1\)。

在 \(\mathfrak p=(X,Y)R\) 处，局部化得到
\[
R_{\mathfrak p}\cong k(Z)[X,Y]_{(X,Y)},\qquad
M_{\mathfrak p}\cong R_{\mathfrak p}\oplus k(Z).
\]
第二个分量成为局部环的剩余域，故同样有非零元素被整个局部极大理想零化，给出深度零。第一个分量是环本身，所以局部模的支集为整个局部谱。这个局部环有素链 \((0)\subsetneq(X)\subsetneq(X,Y)\)，而极大理想由两个元素生成，高度定理给出其维数恰为二。因此在这个支集点，\((S_1)\) 要求深度至少为一，实际却为零。
\end{proof}


\begin{chaptersummary}
极大正则序列的共同长度，可以由固定极大理想生成组的 Koszul 同调最高非零次数确定。正则元素每除去一次，深度和模维数都下降一；深度零则等价于存在被极大理想零化的非零元素。

\ProseParagraphBreak
Ext 把正则序列长度改写为首次非零次数。短正合列于是给出三条深度不等式，各项有限秩的自由消解给出逐次取合冲模的深度下界。比较不同素理想处的深度时，还要计入 \(\dim R/\mathfrak p\)，不能把闭点剩余域直接局部化成另一点的剩余域。

\ProseParagraphBreak
Serre 条件同时约束所有局部化。\((S_1)\) 排除嵌入关联素理想，\((S_2)\) 还要求较高维点处的两步正则性。一般地，\((S_s)\) 模的正则截面满足 \((S_{s-1})\)；特别地，\((S_2)\) 保证正则截面满足 \((S_1)\)。这些条件控制的是局部模结构，不单独等同于约化或正规。
\end{chaptersummary}

\BookExercises

\begin{exercise}\label{b3:ex:16-square-zero-socle}
设 \(k\) 为域，令 \(R=k[U,V]/(U,V)^2\)，其剩余域自然识别为 \(k\)。求 \(\operatorname{Soc}_RR\) 及其作为 \(k\) 向量空间的维数，再求 \(R\) 的 Krull 维数和深度，并计算 \(\dim_k\operatorname{Hom}_R(k,R)\)。
\end{exercise}
\begin{solution}
\(R\) 的 \(k\) 基为 \(1,\bar U,\bar V\)，极大理想 \(\mathfrak m=(\bar U,\bar V)\) 平方为零。一个元素 \(a+b\bar U+c\bar V\) 被 \(\mathfrak m\) 零化当且仅当 \(a=0\)，所以 \(\operatorname{Soc}_RR=\mathfrak m\)，且 \(\dim_k\operatorname{Soc}_RR=2\)。唯一素理想为 \(\mathfrak m\)，故 \(\dim R=0\)；非零基座给出 \(\operatorname{depth}R=0\)。通过 \(f\mapsto f(1)\)，\(\operatorname{Hom}_R(k,R)\) 同构于该二维基座。
\end{solution}

\begin{exercise}\label{b3:ex:16-direct-sum-depth}
设 \((R,\mathfrak m)\) 为交换 Noether 局部环，\(M,N\) 为有限生成的 \(R\) 模。证明
\[
\operatorname{depth}(M\oplus N)=
\min\{\operatorname{depth}M,\operatorname{depth}N\}.
\]
\end{exercise}
\begin{solution}
令 \(k=R/\mathfrak m\)，用一个 \(k\) 的自由消解计算 Ext。Hom 到有限直和逐项分解为两个 Hom 复形，所以
\(\operatorname{Ext}_R^i(k,M\oplus N)\cong\operatorname{Ext}_R^i(k,M)\oplus\operatorname{Ext}_R^i(k,N)\)。
首次非零次数为两者较小者；零模对应永远为零的复形，采用无穷大约定正好保持结论。
\end{solution}

\begin{exercise}\label{b3:ex:16-annihilator-base-ring}
设 \((R,\mathfrak m)\) 为交换 Noether 局部环，\(M\ne0\) 为有限生成的 \(R\) 模，\(I\) 为包含于 \(\operatorname{Ann}_RM\) 的理想。证明作为 \(R\) 模与作为 \(R/I\) 模计算的深度相同，并说明这不表示所有相应的 Ext 模都相同。
\end{exercise}
\begin{solution}
一个 \(\mathfrak m\) 中元素在 \(M\) 及其各次商上的作用只依赖其模 \(I\) 的像。正则序列及不能再延长的条件于是与 \(R/I\) 上完全对应，故深度相同。但对任一域 \(k\)，取 \(R=k[\varepsilon]/(\varepsilon^2)\)、\(I=(\varepsilon)\)、\(M=k\)，在 \(R/I=k\) 上正次数 Ext 为零，在 \(R\) 上则由乘 \(\varepsilon\) 的无限自由消解得到 \(\operatorname{Ext}_R^i(k,k)\cong k\) 对所有 \(i\ge0\) 成立。深度只记录首次非零次数。
\end{solution}

\begin{exercise}\label{b3:ex:16-mixed-components}
设 \(k\) 为域，令 \(A=k[X,Y,Z]_{(X,Y,Z)}/(XY,XZ)\)。证明 \(A\) 的维数为二、深度为一。
\end{exercise}
\begin{solution}
原多项式环中 \((XY,XZ)=(X)\cap(Y,Z)\)，交的两个素理想均为极小素理想：若 \(Xf\in(Y,Z)\)，因 \(X\notin(Y,Z)\) 得 \(f\in(Y,Z)\)，证明了交等式。\(A\) 嵌入两个整环商的直和，故 \(X+Y\) 在 \(A\) 上为非零因子，因为其在两个商中分别为 \(Y,X\)，均非零。两个不可约分量的维数分别为二和一，所以 \(\dim A=2\)。再模掉 \(X+Y\) 得
\[
k[Y,Z]_{(Y,Z)}/(Y^2,YZ),
\]
其中非零元 \(\bar Y\) 被极大理想零化，深度为零。由正则商的深度下降公式，\(\operatorname{depth}A=1\)。
\end{solution}

\begin{exercise}\label{b3:ex:16-nonreduced-positive-depth}
设 \(k\) 为域，令 \(A=k[X,Y,Z]_{(X,Y,Z)}/(X^2,YZ)\)。证明 \(Y+Z\) 在 \(A\) 上正则，求 \(\dim A\)、\(\operatorname{depth}A\)，并计算 \(A/(Y+Z)\) 的长度。说明幂零元和多个不可约分量分别是否阻止了这里的正则性。
\end{exercise}
\begin{solution}
未局部化的商是 \(B=k[Y,Z]/(YZ)\) 上基为 \(1,X\) 的自由模。\(B\) 嵌入 \(k[Y]\oplus k[Z]\)，乘 \(Y+Z\) 在两个分量上分别为乘 \(Y,Z\)，均单射；故它在这个自由 \(B\) 模上也单射，局部化后仍然如此。末端商为
\[
A/(Y+Z)\cong k[X,Y]/(X^2,Y^2),
\]
是剩余域为 \(k\) 的非零 Artin 局部环，基为 \(1,X,Y,XY\)，长度为四。它深度为零，所以正则商的深度下降公式给出 \(\operatorname{depth}A=1\)。去掉幂零元 \(X\) 后得到一维的交叉直线局部环，故 \(\dim A=1\)。\(X\ne0\) 且 \(X^2=0\)，两个极小素理想为 \((X,Y)\)、\((X,Z)\)；这些幂零结构和不同分量均未阻止 \(Y+Z\) 的正则性。
\end{solution}

\begin{exercise}\label{b3:ex:16-dvr-ext}
设 \(R\) 为离散赋值环，统一化元为 \(\pi\)，\(a,b\) 为正整数。计算
\[
\operatorname{Ext}_R^i(R/(\pi^a),R/(\pi^b))\qquad(i=0,1,\ldots),
\]
并求 \(R/(\pi^b)\) 作为 \(R\) 模的深度。
\end{exercise}
\begin{solution}
由 \(0\rightarrow R\xrightarrow{\pi^a}R\rightarrow R/(\pi^a)\rightarrow0\)，Hom 复形是
\[
R/(\pi^b)\xrightarrow{\pi^a}R/(\pi^b).
\]
令 \(c=\min(a,b)\)，其核同构于 \(R/(\pi^c)\)：当 \(a<b\) 时由 \(\pi^{b-a}\) 生成，当 \(a\ge b\) 时为整个模。余核也为 \(R/(\pi^c)\)。因此次数零、一的 Ext 均为此模，其余为零。\(R/(\pi^b)\) 的非零元 \(\pi^{b-1}\) 被极大理想零化，所以深度为零。
\end{solution}

\begin{exercise}\label{b3:ex:16-periodic-ext}
设 \(k\) 为域，令 \(R=k[\varepsilon]/(\varepsilon^2)\)。证明重复乘 \(\varepsilon\) 的序列是剩余域 \(k\) 的自由消解，分别计算 \(\operatorname{Ext}_R^i(k,k)\) 与 \(\operatorname{Ext}_R^i(k,R)\)。
\end{exercise}
\begin{solution}
乘 \(\varepsilon\) 的核与像均为 \((\varepsilon)\)，且增广核也是该理想，故序列正合。Hom 到 \(k\) 后微分全零，所以所有非负次数 Ext 均为 \(k\)。Hom 到 \(R\) 后仍重复乘 \(\varepsilon\)，次数零核为 \((\varepsilon)\cong k\)，正次数核与像相等，所以正次数 Ext 为零。两种系数模的深度都为零，但更高 Ext 行为不同。
\end{solution}

\begin{exercise}\label{b3:ex:16-local-depth-not-monotone}
设 \(k\) 为域，令 \(R=k[X,Y,Z]_{(X,Y,Z)}\)、\(\mathfrak p=(X)R\)。求 \(\operatorname{depth}R\)、\(\operatorname{depth}R_{\mathfrak p}\) 及 \(\dim R/\mathfrak p\)，核对深度的局部化不等式，并与正文中局部化提高深度的例子比较。
\end{exercise}
\begin{solution}
\(X,Y,Z\) 是正则序列，且 \(\dim R=3\)，所以 \(\operatorname{depth}R=3\)。在 \(\mathfrak p\) 处，所有非零的 \(k[Y,Z]\) 元素都成为单位，故
\[
R_{\mathfrak p}\cong k(Y,Z)[X]_{(X)}.
\]
这是维数一的局部整环，\(X\) 正则且最终商为 \(k(Y,Z)\)，故深度一。又 \(R/\mathfrak p\cong k[Y,Z]_{(Y,Z)}\)，维数为二，局部化不等式成为等式 \(3=1+2\)。本例深度降低，正文的直和例子中深度提高；二者说明不能从“局部化”本身预先判断深度数值变化的方向。
\end{solution}

\begin{exercise}\label{b3:ex:16-finite-length-submodule}
设 \(k\) 为域，\(R=k[X,Y]_{(X,Y)}\)、\(M=R/(X^2,XY)\)。求 \(M\) 的最大有限长度子模 \(L\)，并比较 \(M\) 与 \(M/L\) 的深度及维数。
\end{exercise}
\begin{solution}
令 \(L=R\bar X\)。它非零且被 \((X,Y)\) 零化，所以 \(L\cong k\)，长度为一。商 \(M/L\cong R/(X)\cong k[Y]_{(Y)}\) 中乘 \(Y\) 单射，故不含非零有限长度子模：这样的子模会被某次幂 \(Y^n\) 零化，与单射性矛盾。任一有限长度子模 \(L'\subseteq M\) 在 \(M/L\) 中的像仍有限长度，只能为零，因此 \(L'\subseteq L\)。这证明 \(L\) 最大。\(M\) 深度为零，商深度为一；两者维数均为一。除去集中于闭点的有限长度部分改变了深度，却未删去一维的支集分量。
\end{solution}

\begin{exercise}\label{b3:ex:16-ideal-depth}
设 \(k\) 为域，在 \(R=k[X,Y,Z]_{(X,Y,Z)}\) 中，求理想 \(\mathfrak m=(X,Y,Z)\) 和 \(I=(X,Y)\) 作为 \(R\) 模的深度。
\end{exercise}
\begin{solution}
\(R\) 深度为三。由 \(0\rightarrow\mathfrak m\rightarrow R\rightarrow k\rightarrow0\)，中项深度大于末项零，深度引理给出 \(\operatorname{depth}\mathfrak m=1\)。而 \(R/I\cong k[Z]_{(Z)}\) 深度为一，故 \(0\rightarrow I\rightarrow R\rightarrow R/I\rightarrow0\) 给出 \(\operatorname{depth}I=2\)。两者都是整环中的非零理想，支集维数均为三，却有不同深度。
\end{solution}

\begin{exercise}\label{b3:ex:16-three-syzygies}
设 \(k\) 为域，\(R=k[X,Y,Z]_{(X,Y,Z)}\)。用 \(X,Y,Z\) 的 Koszul 消解求剩余域 \(k\) 的第一、第二、第三合冲模的深度。
\end{exercise}
\begin{solution}
第一合冲模为 \(\mathfrak m\)，深度一。第二合冲模来自
\(0\rightarrow\Omega^2k\rightarrow R^3\rightarrow\mathfrak m\rightarrow0\)，因三大于一，深度引理给出深度二。第三合冲模是最左 \(R\) 的单射像，与 \(R\) 同构，深度三。再往前合冲模为零，其深度按约定为无穷大；这一步不能解释为又增加了一个有限的深度。
\end{solution}

\begin{exercise}\label{b3:ex:16-completion-depth}
设 \((R,\mathfrak m)\) 为交换 Noether 局部环，\(M\ne0\) 为有限生成的 \(R\) 模。用极大理想生成组的 Koszul 同调证明
\[
\operatorname{depth}_{\widehat R}\widehat M=\operatorname{depth}_RM.
\]
\end{exercise}
\begin{solution}
第13章给出 \(\widehat M\cong M\otimes_R\widehat R\)，且 \(\widehat R\) 忠实平坦，极大理想为 \(\mathfrak m\widehat R\)。取 \(\mathfrak m=(y_1,\ldots,y_e)\)，同一组像生成完备环的极大理想。平坦张量保持核与像，故
\[
H_i(y;\widehat M)\cong H_i(y;M)\otimes_R\widehat R.
\]
忠实性使两边同时为零。于是 Koszul 最高非零次数不变，生成组长度也不变；由极大正则序列长度公式得到相同深度。
\end{solution}

\begin{exercise}\label{b3:ex:16-s1-torsion-free}
设 \(R\) 为 Noether 整环，\(M\ne0\) 为有限生成的无挠 \(R\) 模。证明 \(M\) 满足 \((S_1)\)，但举例说明它未必满足 \((S_2)\)。
\end{exercise}
\begin{solution}
任意非零 \(m\in M\) 的零化子都是零，故 \(\operatorname{Ass}M=\{0\}\)。其支集为整个素谱，因为 \(\operatorname{Ann}M=0\)；唯一极小素理想也是零，所以 \(M\) 满足 \((S_1)\)。在 \(R=k[X,Y]_{(X,Y)}\) 中取 \(M=(X,Y)\)，该模无挠、维数二、深度一，不满足 \((S_2)\)。
\end{solution}

\begin{exercise}\label{b3:ex:16-s1-not-reduced}
设 \(k\) 为域，令 \(A=k[X,Y,Z]_{(X,Y,Z)}/(X^2,Y^2)\)。证明它是一维非约化局部环，且作为自身模满足每个 Serre 条件 \((S_s)\)（\(s\ge0\)）。
\end{exercise}
\begin{solution}
幂零理想 \((\bar X,\bar Y)\) 的商为 \(k[Z]_{(Z)}\)，故素谱只有极小素理想 \((\bar X,\bar Y)\) 与闭点，两者构成长度一的链。未局部化的商是基为 \(1,X,Y,XY\) 的自由 \(k[Z]\) 模，所以乘 \(Z\) 单射；局部化后商 \(A/ZA=k[X,Y]/(X^2,Y^2)\ne0\)，故闭点深度为一。在极小素点处局部维数为零，条件自动成立；闭点处深度等于维数一，也满足每个 \((S_s)\)。然而 \(\bar X\ne0\)、\(\bar X^2=0\)，所以 \(A\) 不约化。这里条件成立并非只靠零维情形，而是也检查了一维闭点的正则性。
\end{solution}

\begin{exercise}\label{b3:ex:16-maximal-check-insufficient}
设 \(k\) 为域，令 \(R=k[X,Y,Z]_{(X,Y,Z)}\)、\(M=R/(X^2,XY)\)。求闭点及 \(\mathfrak p=(X,Y)R\) 处的深度与模维数，并判断 \(M\) 是否满足 \((S_1)\)。
\end{exercise}
\begin{solution}
未局部化时，这个商是 \(k[X,Y]/(X^2,XY)\) 的多项式扩张，故乘 \(Z\) 单射；闭点处最终商为非零的 \(k[X,Y]_{(X,Y)}/(X^2,XY)\)，其基座含 \(\bar X\)，深度为零。因此闭点深度为一，模维数为二，因为其支集由 \((X)\) 决定。在 \(\mathfrak p\) 处，
\[
M_{\mathfrak p}\cong k(Z)[X,Y]_{(X,Y)}/(X^2,XY).
\]
非零的 \(\bar X\) 被当地极大理想零化，所以深度为零；其支集还有 \((X)\subsetneq(X,Y)\) 的链，维数为一。因此 \((S_1)\) 在这个非闭素点失败，闭点处的正深度没有排除一维嵌入部分。
\end{solution}

\begin{exercise}\label{b3:ex:16-regular-section-embedded}
设 \(k\) 为域，令 \(R=k[X,Y]_{(X,Y)}\)、\(M=(X,Y)\)。证明 \(X\) 是 \(M\) 的非零因子，但 \(M/XM\) 含一个同构于 \(k\) 的子模。解释这与 \(M\) 只满足 \((S_1)\) 的关系。
\end{exercise}
\begin{solution}
\(M\) 嵌入整环 \(R\)，所以乘 \(X\) 单射。\(X\) 的类在 \(M/XM\) 中非零，否则 \(X=Xa\)、\(a\in\mathfrak m\)，约去 \(X\) 得 \(1\in\mathfrak m\)，矛盾。该类被 \(X,Y\) 零化，因为 \(X^2,XY\in XM\)，因此生成一个 \(k\) 子模。模的维数原为二，正则商维数为一，所以该闭点关联素理想是嵌入的，商不满足 \((S_1)\)。正文保证正则商满足 \((S_1)\) 的假设是原模满足 \((S_2)\)，此例恰好不满足。
\end{solution}

\begin{exercise}\label{b3:ex:16-sk-section-general}
设 \(k\) 为域，\(R=k[X,Y]\)、\(M=R\oplus R/(X-1,Y)\)。证明 \(X\) 是 \(M\) 正则元素，\(M/XM\) 满足每个 \((S_s)\)（\(s\ge0\)），但 \(M\) 不满足 \((S_1)\)。说明为何商的 Serre 条件不能直接反推原模的条件。
\end{exercise}
\begin{solution}
\(X\) 在第一直和项上单射，在第二项上为恒等映射，所以乘法单射，且 \(M/XM\cong k[Y]\ne0\)。这个商的支集为 \(V(X)\)，在其极小点处局部模为 \(k(Y)\)、维数零，在每个闭点处为一维局部整环、深度一，所以满足每个 \((S_s)\)。但是在 \(\mathfrak q=(X-1,Y)\) 处，第二直和项是剩余域，给出 \(M_{\mathfrak q}\) 深度零；第一项支集为二维的整个当地素谱，所以 \(\dim M_{\mathfrak q}=2\)，\((S_1)\) 不成立。这里 \(X\notin\mathfrak q\)，取商时整个坏分量被删去；正则商的条件只检查 \(V(X)\) 内的点，不能恢复其外的模结构。
\end{solution}

\begin{exercise}\label{b3:ex:16-sk-direct-sum-failure}
设 \(k\) 为域。在 \(R=k[X,Y]_{(X,Y)}\) 上，说明 \(R\) 与 \(R/(X)\) 都满足 \((S_2)\)，但它们的直和不满足 \((S_2)\)。
\end{exercise}
\begin{solution}
闭点处 \(R\) 深度、维数均为二；高度一处 \(R_{\mathfrak p}\) 是一维局部整环，任意极大理想中的非零元给出深度一；零理想处为域。因此 \(R\) 满足 \((S_2)\)。\(R/(X)\) 的支集只有 \((X)\) 与 \((X,Y)\)，相应局部模分别为域和一维局部整环，深度等于模维数，所以也满足条件。直和在闭点的深度取最小值一，而支集维数取最大值二，故不满足 \((S_2)\)。Serre 条件中的模维数也随直和发生变化。
\end{solution}
